\section{Introducción}
El reconocimiento de patrones y el aprendizaje automático dependen intrínsecamente de la calidad de los datos de entrada \cite{bishop2006,fieguth2022}. En muchos casos empíricos, los datos crudos presentan ruido, escalas dispares y alta dimensionalidad que obstaculizan el rendimiento de los modelos matemáticos y probabilísticos \cite{duda2000}.

Por consiguiente, el preprocesamiento de datos se erige como una fase crítica. Este artículo documenta la construcción de un pipeline estructurado en cuatro etapas principales: procesamiento inicial, normalización, análisis de correlación (reducción de dimensionalidad) y partición de datos. El objetivo principal es transformar un conjunto de datos en bruto en subconjuntos de entrenamiento, validación y prueba óptimos, garantizando la eliminación de redundancias y la escalabilidad del modelo.

El presente estudio evalúa empíricamente diversas estrategias matemáticas sobre dos conjuntos de datos multidimensionales (\textit{Slice} y \textit{Vh}), detallando la toma de decisiones algorítmicas en cada etapa del desarrollo.
